% Source: https://artofproblemsolving.com/community/c3h3164351_2023_emcc_individual_accuracy_test__exeter_math_club_competition
% Stable identity: aops-topic-3164351-q06
The length 5 piece must occupy one complete horizontal row. If it is the middle row, all other pieces must be horizontal in the two remaining rows. Their only partition into row totals 5 is $(1,4)$ and $(2,3)$. Choose which row gets each pair in 2 ways and order each pair in 2 ways, giving 8 tilings. If the length 5 piece is the top or bottom row, the other region is 2-by-5. Pieces 3 and 4 must be horizontal and in different rows. Choose the row containing 4 and the end containing its one-cell gap:4 choices. If piece 2 is horizontal, piece 1 fills that gap and piece 3 can occupy either end of the other row:2 completions. If piece 2 is vertical, it occupies that gap column; the other row then has four consecutive free cells, filled by pieces 3 and 1 in 2 orders:2 more completions. Thus each outer-row placement gives $4(2+2)=16$ tilings. Total $8+16+16=\boxed{40}$.
\begin{asy}
usepackage("amsmath");
import geometry;size(240);
draw((0,0)--(5,0)--(5,3)--(0,3)--cycle);draw((0,1)--(5,1));draw((0,2)--(5,2));draw((4,2)--(4,3));draw((3,1)--(3,2));label("$5$",(2.5,0.5));label("$3$",(1.5,1.5));label("$2$",(4,1.5));label("$4$",(2,2.5));label("$1$",(4.5,2.5));
\end{asy}
